
\documentclass{article}
\usepackage{colortbl}
\usepackage{makecell}
\usepackage{multirow}
\usepackage{supertabular}

\begin{document}

\newcounter{utterance}

\centering \large Interaction Transcript for game `cladder', experiment `full\_v1.5\_default', episode 2646 with refineonly{-}Qwen3.5{-}27B.
\vspace{24pt}

{ \footnotesize  \setcounter{utterance}{1}
\setlength{\tabcolsep}{0pt}
\begin{supertabular}{c@{$\;$}|p{.15\linewidth}@{}p{.15\linewidth}p{.15\linewidth}p{.15\linewidth}p{.15\linewidth}p{.15\linewidth}}
    \# & \multicolumn{2}{c}{Player} && \multicolumn{2}{c}{Game Master} \\
    \hline

    \theutterance \stepcounter{utterance}  
    & & & \multicolumn{4}{p{0.6\linewidth}}{
        \cellcolor[rgb]{0.9,0.9,0.9}{
            \makecell[{{p{\linewidth}}}]{
                \texttt{\tiny{[P1$\langle$GM]}}
                \texttt{You are an expert in causal inference. The following question is not a typical commonsense query, but rather a meticulously designed question created by a professor specializing in causal inference, intended to assess the students' mastery of the course content.} \\
\texttt{Imagine a self{-}contained, hypothetical world with only the following conditions, and without any unmentioned factors or causal relationships: Gender has a direct effect on smoking and lung cancer. Smoking has a direct effect on tar deposit. Tar deposit has a direct effect on lung cancer. Gender is unobserved. For nonsmokers and with no tar deposit, the probability of lung cancer is 79\%. For nonsmokers and with high tar deposit, the probability of lung cancer is 49\%. For smokers and with no tar deposit, the probability of lung cancer is 71\%. For smokers and with high tar deposit, the probability of lung cancer is 40\%. For nonsmokers, the probability of high tar deposit is 79\%. For smokers, the probability of high tar deposit is 40\%. For smokers, would it be more likely to see lung cancer if the person had been a nonsmoker?} \\
\texttt{Start your answer with 'Yes' or 'No', followed by additional reasoning or evidence to support your explanation.} \\
            }
        }
    }
     \\ \\

    \theutterance \stepcounter{utterance}  
    & \multicolumn{4}{p{0.6\linewidth}}{
        \cellcolor[rgb]{0.561,0.737,0.561}{
            \makecell[{{p{\linewidth}}}]{
                \texttt{\tiny{[P1$\rangle$GM]}}
                \texttt{No} \\
\\ 
\texttt{To determine whether a smoker would be more likely to see lung cancer if they had been a nonsmoker, we must compare the observed probability of lung cancer for smokers, \$P(L=1 | S=1)\$, with the counterfactual probability of lung cancer under the intervention of making them nonsmokers, \$P(L=1 | do(S=0))\$.} \\
\\ 
\texttt{**Step 1: Calculate the observed probability of Lung Cancer for Smokers (\$P(L=1 | S=1)\$)**} \\
\texttt{We use the law of total probability, conditioning on Tar (\$T\$):} \\
\texttt{\$\$P(L=1 | S=1) = P(L=1 | S=1, T=0)P(T=0 | S=1) + P(L=1 | S=1, T=1)P(T=1 | S=1)\$\$} \\
\\ 
\texttt{Given data:} \\
\texttt{*   \$P(L=1 | S=1, T=0) = 0.71\$} \\
\texttt{*   \$P(L=1 | S=1, T=1) = 0.40\$} \\
\texttt{*   \$P(T=1 | S=1) = 0.40\$ (Probability of high tar for smokers)} \\
\texttt{*   \$P(T=0 | S=1) = 1 {-} 0.40 = 0.60\$} \\
\\ 
\texttt{Calculation:} \\
\texttt{\$\$P(L=1 | S=1) = (0.71 \textbackslash{}times 0.60) + (0.40 \textbackslash{}times 0.40)\$\$} \\
\texttt{\$\$P(L=1 | S=1) = 0.426 + 0.160 = 0.586\$\$} \\
\\ 
\texttt{**Step 2: Calculate the counterfactual probability of Lung Cancer if the smoker were a nonsmoker (\$P(L=1 | do(S=0))\$)**} \\
\texttt{Under the intervention \$do(S=0)\$, the causal link from Gender to Smoking is severed, and Smoking is fixed to 0. The distribution of Tar (\$T\$) is now determined solely by the intervention \$S=0\$, and the distribution of Lung Cancer (\$L\$) is determined by \$S=0\$ and the resulting \$T\$. The unobserved confounder \$G\$ no longer influences \$S\$ or \$T\$ in this specific calculation because we are forcing \$S=0\$.} \\
\\ 
\texttt{We calculate \$P(L=1 | do(S=0))\$ by conditioning on \$T\$ under the intervention:} \\
\texttt{\$\$P(L=1 | do(S=0)) = P(L=1 | S=0, T=0)P(T=0 | do(S=0)) + P(L=1 | S=0, T=1)P(T=1 | do(S=0))\$\$} \\
\\ 
\texttt{Since \$do(S=0)\$} \\
            }
        }
    }
    & & \\ \\

    \theutterance \stepcounter{utterance}  
    & & & \multicolumn{2}{p{0.3\linewidth}}{
        \cellcolor[rgb]{0.9,0.9,0.9}{
            \makecell[{{p{\linewidth}}}]{
                \texttt{\tiny{[GM$|$GM]}}
                \texttt{No} \\
            }
        }
    }
    & & \\ \\

    \theutterance \stepcounter{utterance}  
    & & & \multicolumn{2}{p{0.3\linewidth}}{
        \cellcolor[rgb]{0.9,0.9,0.9}{
            \makecell[{{p{\linewidth}}}]{
                \texttt{\tiny{[GM$|$GM]}}
                \texttt{no} \\
            }
        }
    }
    & & \\ \\

    \theutterance \stepcounter{utterance}  
    & & & \multicolumn{2}{p{0.3\linewidth}}{
        \cellcolor[rgb]{0.9,0.9,0.9}{
            \makecell[{{p{\linewidth}}}]{
                \texttt{\tiny{[GM$|$GM]}}
                \texttt{game\_result = WIN} \\
            }
        }
    }
    & & \\ \\

\end{supertabular}
}

\end{document}
